% Originally: content/week-11.tex, \section{Differential forms and the exterior derivative} (Corsin Week 11)

\section{Differential forms and the exterior derivative}
\label{sec:differential_forms}

% Source: Corsin Nick/Class Notes/Week 11.pdf, p. 10
\begin{definition}[Differential form]
\label{def:differential_p_form}
A \newterm{differential $p$-form} $\omega$ \germanfor{Differentialform} is defined on a
domain $D \subseteq \mathbb{R}^n$ if, for every $x \in D$, $\omega(x) = \omega_x : (\mathbb{R}^n)^p
\to \mathbb{R}$ is an antisymmetric $p$-linear map. Since $(e_1(x),\ldots,e_n(x))$ is a basis
of $\mathbb{R}^n$ for every $x \in D$, \cref{prop:k_form_basis_expansion} lets us write
\[\omega(x) = \sum_{1\leq i_1<\cdots<i_p\leq n} \omega_{i_1\cdots i_p}(x)\,e_{i_1}^*(x)\wedge
  \cdots\wedge e_{i_p}^*(x)\]
for functions $\omega_{i_1\cdots i_p} : \mathbb{R}^n \to \mathbb{R}$.
\end{definition}

% Source: Corsin Nick/Class Notes/Week 11.pdf, p. 11
\begin{notation}
We abbreviate $e_{i_1}^*\wedge\cdots\wedge e_{i_p}^* =: e_I^*$ for $I = (i_1,\ldots,i_p)$, and
write $\Omega^p(U)$ for the set of $p$-forms on $U \subseteq \mathbb{R}^n$. By convention,
$\Omega^0(U) = C^\infty(U)$.
\end{notation}

\begin{example}
Let $U \subseteq \mathbb{R}^n$.
\begin{enumerate}[label=\textbf{(\alph*)}]
  \item Any $f \in C^\infty(U)$ is a $0$-form.
  \item A vector field $F : U \to \mathbb{R}^n$ has an associated $1$-form $\omega_F(x) : v
  \mapsto \langle F(x),v\rangle$, that is,
  \[\omega_F(x)(v) = \sum_{i=1}^n F^i(x)\,v^i = \sum_{i=1}^n F^i(x)\,e_i^*(x)(v)
    \qquad\Longrightarrow\qquad \omega_F = \sum_{i=1}^n F^i\,e_i^*.\]
% Source: Corsin Nick/Class Notes/Week 11.pdf, p. 12
  \item $F : U \to \mathbb{R}^n$ also naturally induces an $(n-1)$-form
  \[\omega_F^{n-1}(x)(v_1,\ldots,v_{n-1}) := \det\big(F(x)\mid v_1\mid\cdots\mid v_{n-1}\big).\]
  Developing this determinant along the first column,
  \begin{align}
  \omega_F^{n-1}(x)(v_1,\ldots,v_{n-1}) &= \sum_{i=1}^n (-1)^{i-1}F^i(x)\det\!\left(
    \begin{smallmatrix} v_1^1 & \cdots & v_{n-1}^1 \\ \vdots & & \vdots \\ v_1^{i-1} & \cdots &
    v_{n-1}^{i-1} \\ v_1^{i+1} & \cdots & v_{n-1}^{i+1} \\ \vdots & & \vdots \\ v_1^n &
    \cdots & v_{n-1}^n \end{smallmatrix}\right) \nonumber \\
    &= \sum_{i=1}^n (-1)^{i-1}F^i(x)\,e_1^*\wedge\cdots\wedge\widehat{e_i^*}\wedge\cdots\wedge
    e_n^*(v_1,\ldots,v_{n-1}), \nonumber
  \end{align}
  where $\widehat{e_i^*}$ denotes an omitted factor, so that $\omega_F^{n-1} =
  \sum_{i=1}^n (-1)^{i-1}F^i\,e_1^*\wedge\cdots\wedge\widehat{e_i^*}\wedge\cdots\wedge e_n^*$.
  \item In particular, for $F : \mathbb{R}^3 \to \mathbb{R}^3$,
  \[\omega_F^2(x)(v_1,v_2) = \langle F(x),v_1\times v_2\rangle = \det\big(F(x)\mid v_1\mid
    v_2\big).\]
\end{enumerate}
\end{example}

\begin{ainote}
Item \textbf{(d)} identifies $\omega_F^2$ with the flux integrand from
\cref{sec:flux_divergence_theorem}: the flux $\int_S \langle F,\nu\rangle\,d\vol_2$ over a parametrized surface is
exactly the integral, in the sense of the motivating change-of-variables calculation above, of
the $2$-form $\omega_F^2$.
\end{ainote}

% Source: Corsin Nick/Class Notes/Week 11.pdf, p. 13
\begin{theorem}[Exterior derivative]
\label{thm:exterior_derivative}
There exists a unique linear operator $d : \Omega^k(U) \to \Omega^{k+1}(U)$ satisfying
\begin{enumerate}[label=\textbf{(\alph*)}]
  \item for $f \in C^\infty(U) = \Omega^0(U)$: $df = Df$, the ordinary differential;
  \item for $\omega \in \Omega^k(U)$, $\theta \in \Omega^s(U)$: $d(\omega\wedge\theta) =
  d\omega\wedge\theta + (-1)^k\,\omega\wedge d\theta$;
  \item $d\circ d\,\omega = d^2\omega = 0$.
\end{enumerate}
\end{theorem}

\begin{ainote}
Corsin states this as a fact seen in the lecture, without proof, and moves directly to its
consequences. The proof can be found in the official script
(Analysis\_II\_Script\_v1.pdf, §14.3, pp.\ 136--138), where the operator is constructed
explicitly in coordinates and properties \textbf{(a)}--\textbf{(c)} are verified.
\end{ainote}

Denote by $dx^i$ the differential of $x \mapsto x^i$ in a moving frame $(e_1(x),\ldots,e_n(x))$,
that is, $dx^i = e_i^*(x)$. Every $p$-form $\omega$ on $U \subseteq \mathbb{R}^n$ then takes
the shape
\[\omega = \sum_{1\leq i_1<\cdots<i_p\leq n} \omega_{i_1\cdots i_p}\,dx^{i_1}\wedge\cdots\wedge
  dx^{i_p} = \sum_I \omega_I\,dx^I, \qquad \omega_I \in C^\infty(U),\]
with exterior derivative
\[d\omega = \sum_I d\omega_I \wedge dx^I.\]

% Source: Corsin Nick/Class Notes/Week 11.pdf, pp. 14--15
\begin{example}[The exterior derivative recovers the gradient and the curl]
\label{ex:exterior_derivative_gradient_curl}
\begin{enumerate}[label=\textbf{(\alph*)}]
  \item For $f \in C^\infty(\mathbb{R}^n) = \Omega^0(\mathbb{R}^n)$,
  \[df_x(v) = \left(\frac{\partial f}{\partial x^1}\ \Big|\ \cdots\ \Big|\ \frac{\partial f}
    {\partial x^n}\right)\!\begin{pmatrix}v^1\\ \vdots \\ v^n\end{pmatrix}
    = \frac{\partial f}{\partial x^1}\,dx^1(v) + \cdots + \frac{\partial f}{\partial x^n}\,
    dx^n(v),\]
  so $df = \frac{\partial f}{\partial x^1}\,dx^1 + \cdots + \frac{\partial f}{\partial x^n}\,
  dx^n = \langle \nabla f,\cdot\rangle$: the exterior derivative of a $0$-form is exactly the
  $1$-form associated (as in the example above) to the gradient vector field $\nabla f$. This
  is the shape that makes Stokes' theorem specialize to the FTC, as in
  \cref{sec:gauss_intuition}.
  \item Define a $1$-form on $U \subseteq \mathbb{R}^3$ by $\omega(x,y,z) = \langle
  F(x,y,z),\cdot\rangle = F^1\,dx+F^2\,dy+F^3\,dz$ for a vector field $F$. By property
  \textbf{(b)} of
  \cref{thm:exterior_derivative} and $d(dx^i)=0$ (since $dx^i = d(x^i)$ and $d^2=0$),
  \[
  d\omega = dF^1\wedge dx + dF^2\wedge dy + dF^3\wedge dz,
  \]
  and expanding each $dF^i = \partial_x F^i\,dx+\partial_y F^i\,dy+\partial_z F^i\,dz$, only
  the cross terms survive the wedge (since $dx^i\wedge dx^i=0$ by antisymmetry):
  \[
  d\omega = \Big(\frac{\partial F^3}{\partial y}-\frac{\partial F^2}{\partial z}\Big)dy\wedge
    dz + \Big(\frac{\partial F^1}{\partial z}-\frac{\partial F^3}{\partial x}\Big)dz\wedge dx
    + \Big(\frac{\partial F^2}{\partial x}-\frac{\partial F^1}{\partial y}\Big)dx\wedge dy.
  \]
  Comparing with example item \textbf{(d)} in \cref{sec:differential_forms}, this is exactly
  $\omega_{\nabla\times F}^2 = \langle\nabla\times F,\cdot\times\cdot\rangle = \det(\nabla
  \times F\mid\cdot\mid\cdot)$: the exterior derivative of the $1$-form associated to $F$ is
  the $2$-form associated to $\curl F$.

  % Generator: Claude Sonnet 5 (High) --- new content: the analogous fact for the divergence,
  % completing the grad/curl/div trio as the 0-, 1- and (n-1)-form faces of $d$, using the
  % $(n-1)$-form $\omega_F^{n-1}$ already introduced in \cref{sec:differential_forms} (example
  % item (c)) rather than redefining it. Used below in
  % \cref{rem:one_theorem_four_costumes} to place the divergence theorem alongside Green,
  % classical Stokes and the FTC as instances of the single general Stokes theorem.
  \item Recall the $(n-1)$-form $\omega_F^{n-1} =
  \sum_{i=1}^n(-1)^{i-1}F^i\,dx^1\wedge\cdots\wedge\widehat{dx^i}\wedge\cdots\wedge dx^n$ from
  \cref{sec:differential_forms} (example item \textbf{(c)}), with $\widehat{dx^i}$ denoting an
  omitted factor. In $d(F^i\,dx^1\wedge\cdots\widehat{dx^i}\cdots\wedge dx^n) =
  dF^i\wedge dx^1\wedge\cdots\widehat{dx^i}\cdots\wedge dx^n$, only the $\partial_iF^i\,dx^i$
  summand of $dF^i = \sum_j\partial_jF^i\,dx^j$ survives, since every $j\neq i$ repeats a factor
  already present and so kills the wedge by antisymmetry. Therefore
  \[d\omega_F^{n-1} = \sum_{i=1}^n(-1)^{i-1}\partial_iF^i\;dx^i\wedge dx^1\wedge\cdots\wedge
    \widehat{dx^i}\wedge\cdots\wedge dx^n.\]
  Moving $dx^i$ past the $i-1$ factors $dx^1,\dots,dx^{i-1}$ ahead of it, into its natural
  position, costs another sign of $(-1)^{i-1}$, cancelling the one already there:
  \begin{align*}
    dx^i\wedge dx^1\wedge\cdots\wedge\widehat{dx^i}\wedge\cdots\wedge dx^n
      &= (-1)^{i-1}\,dx^1\wedge\cdots\wedge dx^n \\
    \Longrightarrow\quad
      d\omega_F^{n-1} &= \sum_{i=1}^n\partial_iF^i\;dx^1\wedge\cdots\wedge dx^n
      = (\divg F)\,dx^1\wedge\cdots\wedge dx^n.
  \end{align*}
  So, the exterior derivative of $\omega_F^{n-1}$ is $\divg F$ times the volume form. That
  completes the pattern of items \textbf{(a)} and \textbf{(b)} above: $\grad$, $\curl$ and
  $\divg$ are exactly the $0$-, $1$- and $(n-1)$-form faces of the single operator $d$.

% Source: Jérôme Paschoud/Notizen/Woche 12.1 Differentialformen.pdf, p. 6
  \item Let $\omega = x_3 \,dx_3 + x_1 x_2 \,dx_2$ be a $1$-form on $\mathbb{R}^3$. Let us compute its exterior derivative using the rules:
  \begin{align*}
  d\omega &= d(x_3) \wedge dx_3 + d(x_1 x_2) \wedge dx_2 \\
  &= dx_3 \wedge dx_3 + (x_2 \,dx_1 + x_1 \,dx_2) \wedge dx_2 \\
  &= 0 + x_2 \,dx_1 \wedge dx_2 + x_1 \underbrace{dx_2 \wedge dx_2}_{=0} \\
  &= x_2 \,dx_1 \wedge dx_2.
  \end{align*}
\end{enumerate}
\end{example}

\begin{ainote}
Corsin writes the middle expansion of $d\omega$ out in full, all nine terms, one for each pair
$(dF^i,dx^j)$, before collecting the six that survive into the three coefficient pairs above. It
is condensed here to the surviving terms directly. Nothing is lost, since $dx\wedge
dx=dy\wedge dy=dz\wedge dz=0$ kills exactly the three diagonal terms, and antisymmetry
($dy\wedge dx=-dx\wedge dy$, and so on) pairs up the rest.
\end{ainote}

% Generator: Claude Opus 5 (High) --- FIG-CH24-03. The chapter opens by promising that grad,
% curl and div are one operator satisfying one identity. Items (a)-(c) of the example above prove
% it, one column at a time, and the result was never assembled anywhere. This is that assembly.
% Every vertical arrow is an identification established earlier in this chapter, and each is
% labelled with the place it was made, so the diagram is checkable rather than decorative.
\begin{center}
\begin{tikzcd}[row sep=3.4em, column sep=3.4em]
  \Omega^0(U) \arrow[r, "d"]
    & \Omega^1(U) \arrow[r, "d"]
    & \Omega^2(U) \arrow[r, "d"]
    & \Omega^3(U) \\
  C^\infty(U) \arrow[r, "\grad"'] \arrow[u, "="]
    & C^\infty(U,\mathbb{R}^3) \arrow[r, "\curl"'] \arrow[u, "F \mapsto \omega_F"]
    & C^\infty(U,\mathbb{R}^3) \arrow[r, "\divg"'] \arrow[u, "F \mapsto \omega_F^2"]
    & C^\infty(U) \arrow[u, "g \mapsto g\,dx\wedge dy\wedge dz"']
\end{tikzcd}
\end{center}

\begin{importantremark}[One operator, one identity]
\label{rem:one_theorem_four_costumes_forms}
For $U \subseteq \mathbb{R}^3$ the diagram above commutes, and each square is one of the three
computations just performed: the left square is item \textbf{(a)}, the middle is item
\textbf{(b)}, the right is item \textbf{(c)}. Read downward it says that the classical vector
operators are $d$ in disguise, and that the disguise is nothing but the choice of which object to
call a vector field.

The payoff is the bottom row's two mysteries collapsing into one line of the top row. Property
\textbf{(c)} of \cref{thm:exterior_derivative} is $d \circ d = 0$, a single statement about a
single
operator, and reading it through the two identifications gives
\[\curl\grad = 0 \qquad\text{and}\qquad \divg\curl = 0.\]
These are the same fact twice, seen from two rungs of the same ladder. That is the whole reason
for the machinery: in $\mathbb{R}^3$ there happen to be two identities because there happen to be
two interior rungs, and in $\mathbb{R}^n$ there are $n-1$ of them, none of which needs a separate
proof.

Two cautions the picture invites. The vertical arrows are \emph{not} canonical: $F \mapsto
\omega_F$ uses the inner product and $F \mapsto \omega_F^2$ uses the determinant, so both depend
on $\mathbb{R}^3$ carrying a metric and an orientation, which is exactly what a general manifold
does not hand you for free. And the ladder is this short only because $n = 3$: $\Omega^4(U)$ on
$\mathbb{R}^3$ is zero, since there is no fourth coordinate to build a shadow on
(\cref{rem:forms_measure_shadows}), which is why $\divg$ is the last operator in the row and not
because anyone chose to stop there.
\end{importantremark}

\begin{aiexercise}[The angle form is closed]
% Generator: Claude Opus 5 (High)
\label{ex:ai_angle_form_closed}
On $U := \mathbb{R}^2\setminus\{0\}$ define the $1$-form
\[\omega := \frac{-y\,dx + x\,dy}{x^2+y^2}.\]
\begin{enumerate}[label=\textbf{(\alph*)}]
  \item Show that $d\omega = 0$.
  \item By part \textbf{(a)} and \cref{thm:exterior_derivative}, one might hope that
    $\omega = df$ for some
    $f \in C^\infty(U)$. Compute $\int_\gamma \omega$ along the unit circle
    $\gamma(t) := (\cos t, \sin t)$, $t \in [0,2\pi]$, and explain what this rules out.
    \emph{Hint:} \cref{prop:work_difference_potential}, read through the identification
    $F \leftrightarrow \omega_F$ of the diagram above.
\end{enumerate}
\end{aiexercise}

\begin{aiexercise}[Exterior derivatives by hand]
% Generator: Claude Opus 5 (High)
\label{ex:ai_exterior_derivative_drill}
Compute $d\omega$ for each of the following, on the largest domain where the expression makes
sense.
\begin{enumerate}[label=\textbf{(\alph*)}]
  \item $\omega = xy\,dx + yz\,dy + zx\,dz$ on $\mathbb{R}^3$. Which vector field's curl have you
    just computed?
  \item $\omega = x_1x_2x_3\,dx_2\wedge dx_3$ on $\mathbb{R}^3$.
  \item $\omega = f\,dg$, where $f,g \in C^\infty(\mathbb{R}^n)$. Express the answer using
    $\wedge$ and deduce that $d(f\,dg + g\,df) = 0$.
\end{enumerate}
\end{aiexercise}

% Source: DiffComp.pdf (Prof. Joaquim Serra, "Some differential form computations 1",
% Analysis II: Several Variables, ETH Zürich FS 2026, assignment 13 May 2026, not graded), pp. 1--2
% Extractor: Claude Opus 5 (High)
% This is our OWN lecturer's example sheet for the current course, not an old exam, and it was
% never cited anywhere in the repository before 2026-08-11. The six parts are Serra's, verbatim in
% content; only the numbering changes (his 1.1--1.6 become our (a)--(f), per the sub-part rule).
\begin{exercise}[Six wedge products $df \wedge \lambda$]
\label{ex:serra_df_wedge_lambda}
All forms below live on $\mathbb{R}^3$, and in every case $\lambda$ is a $1$-form. Compute
\[\omega := df \wedge \lambda\]
in simplified form. In parts \textbf{(a)}--\textbf{(c)} the symbol $f$ denotes a function, that is
a $0$-form, and the answer should be presented as
\[\omega = F\,dx\wedge dy + G\,dx\wedge dz + H\,dy\wedge dz;\]
in parts \textbf{(d)}--\textbf{(f)} the symbol $f$ denotes a $1$-form instead, and the answer
should be presented as
\[\omega = F\,dx\wedge dy\wedge dz.\]
\begin{enumerate}[label=\textbf{(\alph*)}]
  \item $f(x,y) = x^2y + \sin y$, \qquad $\lambda = x\,dx + y\,dy$.
  \item $f(x,y,z) = xyz$, \qquad $\lambda = x\,dy + y\,dz + z\,dx$.
  \item $f(x,y,z) = e^xy + z^2$, \qquad $\lambda = (x+y)\,dx + z\,dy + xz\,dz$.
  \item $f = x^2\,dy + y^2\,dz$, \qquad $\lambda = x\,dx + z\,dy$.
  \item $f = xy\,dx + z\,dy + x^2z\,dz$, \qquad $\lambda = y\,dx + x\,dy + z\,dz$.
  \item $f = e^y\,dx + xz\,dy + y^2z\,dz$, \qquad $\lambda = dx + y\,dy + z\,dz$.
\end{enumerate}
\exinfo{This exercise is the whole of \emph{Some differential form computations 1}, an ungraded
example sheet issued by Prof.\ Joaquim Serra on 13 May 2026 for Analysis II, Spring Semester 2026.
The sheet prints full worked solutions, which are the basis of the solutions given here.}
\end{exercise}

% Feedback: Gemini 3.1 Pro (2026-08-11) --- called a very strong intervention: warning the reader
% that Serra's $f$ is a $0$-form in the first half and a $1$-form in the second heads off immense
% confusion. Keep this note; do not "tidy" the exercise by renaming Serra's letters instead.
\begin{ainote}
Serra uses the letter $f$ for a $0$-form in \textbf{(a)}--\textbf{(c)} and for a $1$-form in
\textbf{(d)}--\textbf{(f)}, and the sheet says so explicitly rather than changing the name. Keep
the distinction in view while computing, because $df$ means two different things across the two
halves: the differential of a function in the first, and the exterior derivative of a $1$-form in
the second (\cref{thm:exterior_derivative}). The degree bookkeeping is what makes the answers land
in different spaces: $df\wedge\lambda$ is a $2$-form in \textbf{(a)}--\textbf{(c)} and a $3$-form
in \textbf{(d)}--\textbf{(f)}, which is why the requested output shapes differ.
\end{ainote}

% Supplement: Sascha Brack/Class Notes/Week_11_Notes_Friday.pdf, pp. 2--4
